\section*{Capítulo \ref{ch:b1:organomagnesium} --- Os reagentes organomagnesianos}

\begin{solution}{exo:b1:organomagnesium:1}
\ce{CH3CH2Br + Mg -> CH3CH2MgBr}; \ce{C6H5Br + Mg -> C6H5MgBr} (em éter
seco). C--Br: Br (2.96) mais eletronegativo que C (2.55), carbono
positivo. C--Mg: Mg (1.31) menos eletronegativo, carbono negativo.
\end{solution}

\begin{solution}{exo:b1:organomagnesium:2}
Água: \ce{CH3MgI + H2O -> CH4 + Mg(OH)I}.
Etanol: \ce{CH3MgI + CH3CH2OH -> CH4 + CH3CH2OMgI}.
Ácido etanoico: \ce{CH3MgI + CH3COOH -> CH4 + CH3COOMgI}.
Óxido de deutério: \ce{CH3MgI + D2O -> CH3D + Mg(OD)I}.
\end{solution}

\begin{solution}{exo:b1:organomagnesium:3}
Metanal: propan-1-ol (primário). Etanal: butan-2-ol (secundário).
Propanona: 2-metilbutan-2-ol (terciário). Benzaldeído:
1-fenilpropan-1-ol (secundário). Dióxido de carbono: ácido propanoico.
\end{solution}

\begin{solution}{exo:b1:organomagnesium:4}
Seu grupo OH é ácido o bastante para destruir o reagente assim que ele
se forma: a molécula protonaria a si mesma (ou suas vizinhas), dando o
alcóxido de magnésio do butan-1-ol, \ce{CH3(CH2)3OMgBr}, e não um
reagente utilizável. O OH precisa antes ser protegido
(o \cref{ch:b1:acetals-protection} mostra a ideia para um aldeído).
\end{solution}

\begin{solution}{exo:b1:organomagnesium:5}
O par C--Mg ataca o carbono carbonílico do etanal, enquanto o par $\pi$ da
C=O vai para O: \ce{CH3CH(CH3)O^-} com \ce{MgBr+}; depois \ce{H3O+} cede um
próton ao oxigênio: propan-2-ol. O propan-2-ol leva dois grupos metila
no carbono carbinólico: ele é \omterm{def:b1:stereochemistry-in-depth:stereocentre}{aquiral}. (Com outro R, por exemplo o
brometo de etilmagnésio com o etanal, o produto, o butan-2-ol, é \omterm{def:b1:stereochemistry-in-depth:stereocentre}{quiral},
mas racêmico: a carbonila plana é atacada igualmente pelas duas faces.)
\end{solution}

\begin{solution}{exo:b1:organomagnesium:6}
O carbono carbinólico leva um grupo metila e dois grupos etila: brometo
de etilmagnésio com butan-2-ona, ou brometo de metilmagnésio com
pentan-3-ona.
\end{solution}

\begin{solution}{exo:b1:organomagnesium:7}
Prepare \ce{(CH3)2CHMgBr} a partir do 2-bromopropano e de magnésio em
éter seco, derrame-o sobre \ce{CO2} sólido e depois acidifique:
\ce{(CH3)2CHCOOH}. $n = 12.3/123.0 = 0.100$ mol;
$0.100 \times 0.70 \times 88.0 = \qty{6.2}{g}$.
\end{solution}

\begin{solution}{exo:b1:organomagnesium:8}
$0.18/18.0 = \qty{0.010}{mol}$ de água destrói \qty{0.010}{mol} de
reagente: \qty{20}{\percent}; forma-se benzeno, \ce{C6H5MgBr + H2O -> C6H6 +
Mg(OH)Br}.
\end{solution}

\begin{solution}{exo:b1:organomagnesium:9}
O reagente, um \omterm{def:b1:electronic-effects:nucleophile}{nucleófilo}, reage com o bromobenzeno ainda presente: no
total, \ce{C6H5MgBr + C6H5Br -> C6H5-C6H5 + MgBr2}. Adicionar o brometo
lentamente a um excesso de magnésio mantém sua concentração baixa: ele
encontra metal novo, e não o reagente já formado.
\end{solution}

\begin{solution}{exo:b1:organomagnesium:10}
1-Feniletan-1-ol: \ce{CH3MgBr} (a partir do bromometano) com benzaldeído,
ou \ce{C6H5MgBr} (a partir do bromobenzeno) com etanal.
2-Fenilpropan-2-ol: \ce{C6H5MgBr} com propanona, ou \ce{CH3MgBr} com
1-feniletan-1-ona.
\end{solution}

\begin{solution}{exo:b1:organomagnesium:11}
A cetona é adicionada ao reagente para que o reagente, preparado no
balão seco, nunca saia dele e para que o calor liberado por cada porção
seja absorvido pelo éter em ebulição. Só a água daria o sal básico
gelatinoso \ce{Mg(OH)Br}, que aprisiona o produto; o ácido diluído frio
o dissolve como \ce{Mg^2+}, e o frio limita as reações secundárias do
álcool terciário em meio ácido.
\end{solution}

\begin{solution}{exo:b1:organomagnesium:12}
Cada mol de reagente dá com a água um mol de sal básico de magnésio,
\ce{RMgBr + H2O -> RH + Mg(OH)Br}, cujo \ce{OH-} é titulado pelo ácido:
a quantidade de reagente é igual à de ácido usado, \qty{0.088}{mol}, um
\omterm{def:b1:extent-q-and-k:fractional-extent}{rendimento} de
\qty{88}{\percent}.
\end{solution}

\begin{solution}{pb:b1:organomagnesium:1}
\textbf{1.} \ce{C6H5Br + Mg -> C6H5MgBr}.
\textbf{2.} Mg $0.535/24.3 = \qty{0.0220}{mol}$; \ce{C6H5Br} $3.14/156.9 =
\qty{0.0200}{mol}$: magnésio em ligeiro excesso.
\textbf{3.} A água destrói o reagente: \ce{C6H5MgBr + H2O -> C6H6 +
Mg(OH)Br}.
\textbf{4.} Ele impede que o vapor d'água do ar entre pelo condensador.
\textbf{5.} O éter cede \omterm{def:b1:lewis-resonance-vsepr:lewis}{pares não ligantes} ao magnésio (\omterm{def:b1:lewis-resonance-vsepr:lewis-acid}{base de Lewis}):
sem ele, o reagente não se forma nem permanece em solução.
\textbf{6.} \ce{C6H5MgBr + C6H5Br -> C6H5-C6H5 + MgBr2}; bifenila
$0.15/154.0 = \qty{9.7e-4}{mol}$, que consumiu \qty{9.7e-4}{mol} de
bromobenzeno e o mesmo tanto de reagente: \qty{1.9e-3}{mol} de
bromobenzeno perdidos no total.
\textbf{7.} $3.28/182.0 = \qty{0.0180}{mol}$.
\textbf{8.} O par C--Mg ataca o carbono carbonílico, o par $\pi$ vai para
o oxigênio: \ce{(C6H5)3C-O^-} \ce{MgBr+}.
\textbf{9.} Reagente disponível no máximo $0.0200 - 0.0019 = \qty{0.0181}{mol}$;
benzofenona \qty{0.0180}{mol}: a benzofenona é limitante, por pouco.
\textbf{10.} Não: o carbono carbinólico leva três grupos fenila idênticos.
\textbf{11.} O reagente fica em seu balão seco, e o calor de cada porção
de cetona é absorvido pelo éter em ebulição.
\textbf{12.} Destruiria parte do reagente (benzeno) e deixaria
benzofenona sem reagir.
\textbf{13.} \ce{(C6H5)3COMgBr + H3O+ -> (C6H5)3COH + Mg^2+ + Br- + H2O}.
\textbf{14.} O ácido dissolve os sais básicos de magnésio que só a água
precipitaria; o frio limita a formação do \omterm{def:b1:electronic-effects:intermediates}{carbocátion} da questão 18.
\textbf{15.} O trifenilmetanol e a bifenila na \omterm{def:b1:extent-q-and-k:system}{fase} etérea; os sais de
magnésio na \omterm{def:b1:extent-q-and-k:system}{fase} aquosa.
\textbf{16.} Lavando (ou triturando) o sólido com éter de petróleo frio:
a bifenila se dissolve, o trifenilmetanol fica.
\textbf{17.} Um ponto de fusão nítido e igual ao tabelado (ou uma
única mancha em cromatografia em camada delgada).
\textbf{18.} \ce{(C6H5)3COH + H+ -> (C6H5)3C+ + H2O}: um \omterm{def:b1:electronic-effects:intermediates}{carbocátion}
terciário cuja carga se espalha por três anéis benzênicos por
ressonância, muito estável e colorido por causa dessa conjugação
estendida.
\textbf{19.} $3.75/260.0 = \qty{0.0144}{mol}$.
\textbf{20.} $0.0144/0.0180 = \qty{80}{\percent}$.
\textbf{21.} Perdas nas transferências e na recristalização; reagente
destruído por traços de umidade; reação incompleta.
\textbf{22.} O ácido benzoico, \ce{C6H5COOH}, após acidificação.
\textbf{23.} $0.0181 \times 122.0 = \qty{2.21}{g}$.
\textbf{24.} \ce{C6H5MgBr + CO2 -> C6H5COOMgBr}, depois
\ce{C6H5COOMgBr + H3O+ -> C6H5COOH + Mg^2+ + Br- + H2O}.
\textbf{25.} \textbf{\qty{80}{\percent}}.
\end{solution}
